\documentclass[9pt,a4paper]{extarticle}
\usepackage{parskip}

% ======================
% Packages
% ======================
\usepackage[utf8]{inputenc}
\usepackage[T1]{fontenc}
\usepackage[brazil]{babel}
\usepackage{amsmath, amssymb, amsthm}
\usepackage{geometry}
\usepackage{listings}
\usepackage{xcolor}
\usepackage{hyperref}
\usepackage{booktabs}
\usepackage{array}
\usepackage{tabularx}

\geometry{margin=2cm}

% ======================
% Listings (Python)
% ======================
\definecolor{codebg}{rgb}{0.95,0.95,0.95}
\lstset{
  language=Python,
  backgroundcolor=\color{codebg},
  basicstyle=\ttfamily\footnotesize,
  keywordstyle=\color{blue},
  commentstyle=\color{gray},
  stringstyle=\color{red!70!black},
  showstringspaces=false,
  frame=single,
  breaklines=true,
  inputencoding=utf8
}

% ======================
% Title
% ======================
\begin{document}

\title{Roteiro de Aula --- Classifica\c c\~ao de Matrizes via Forma Quadr\'atica}
\author{Disciplina: PCC179/BCC405 --- Otimiza\c c\~ao N\~ao Linear}
\date{\today}
\maketitle

\section{Motiva\c c\~ao}
Em otimiza\c c\~ao n\~ao linear, a \textbf{Hessiana} $\nabla^2 f(x^\ast)$ em um ponto cr\'itico $x^\ast$ determina a natureza do ponto:
\begin{itemize}
  \item positiva definida $\Rightarrow$ m\'inimo estrito;
  \item negativa definida $\Rightarrow$ m\'aximo estrito;
  \item indefinida $\Rightarrow$ ponto de \textit{sela};
  \item semidefinida $\Rightarrow$ casos degenerados (an\'alise adicional).
\end{itemize}
Isso motiva a \textbf{classifica\c c\~ao de matrizes sim\'etricas} reais $A \in \mathbb{R}^{n\times n}$ via o sinal da forma quadr\'atica $q_A(x)=x^\top A x$.

\section{Forma Quadr\'atica e Defini\c c\~oes}
Seja $A=A^\top \in \mathbb{R}^{n\times n}$. Definimos:
\[
q_A(x) = x^\top A x, \quad x\in\mathbb{R}^n.
\]
\begin{itemize}
  \item \textbf{Positiva definida (PD):} $x^\top A x>0$ para todo $x\neq 0$.
  \item \textbf{Positiva semidefinida (PSD):} $x^\top A x\ge 0$ para todo $x$.
  \item \textbf{Negativa definida (ND):} $x^\top A x<0$ para todo $x\neq 0$.
  \item \textbf{Negativa semidefinida (NSD):} $x^\top A x\le 0$ para todo $x$.
  \item \textbf{Indefinida:} existem $u,v$ tais que $u^\top A u>0$ e $v^\top A v<0$.
\end{itemize}

\section{Crit\'erios de Classifica\c c\~ao}
\subsection{Crit\'erio dos Autovalores}

\subsection*{1) Conceito e C\'alculo}
Para uma matriz sim\'etrica $A \in \mathbb{R}^{n\times n}$, os \textbf{autovalores} $\lambda_i$ e \textbf{autovetores} $v_i$ s\~ao definidos por:
\[
A v_i = \lambda_i v_i, \quad i=1,\dots,n.
\]
Os autovalores s\~ao obtidos resolvendo o \textbf{polin\^omio caracter\'istico}:
\[
\det(A - \lambda I) = 0.
\]
Cada raiz $\lambda_i$ representa um fator de escala associado ao autovetor $v_i$, indicando quanto a transforma\c c\~ao linear $A$ alonga ou contrai essa dire\c c\~ao.

\paragraph{M\'etodos de c\'alculo.}
Para matrizes pequenas, pode-se resolver simbolicamente o determinante.  
Em pr\'atica, algoritmos num\'ericos (como os baseados em decomposi\c c\~oes QR ou Jacobi) s\~ao usados, e bibliotecas como o \texttt{NumPy} oferecem fun\c c\~oes eficientes (ex.: \texttt{np.linalg.eigvalsh} para matrizes sim\'etricas).  

\begin{itemize}
  \item Se $A$ \'e real e sim\'etrica, todos os autovalores $\lambda_i$ s\~ao reais.
  \item Os autovetores formam uma base ortogonal do espa\c co.
\end{itemize}

\subsection*{2) Exemplo Num\'erico}

Considere a matriz:
\[
A = 
\begin{bmatrix}
2 & -1 & 0 \\
-1 & 2 & -1 \\
0 & -1 & 2
\end{bmatrix}.
\]
O polin\^omio caracter\'istico \'e:
\[
\det(A - \lambda I) =
\begin{vmatrix}
2-\lambda & -1 & 0 \\
-1 & 2-\lambda & -1 \\
0 & -1 & 2-\lambda
\end{vmatrix}
= (2-\lambda)\big[(2-\lambda)^2 - 1\big] - (-1)[(-1)(2-\lambda)].
\]
Simplificando:
\[
\det(A - \lambda I) = (2-\lambda)^3 - 2(2-\lambda) = (2-\lambda)\big[(2-\lambda)^2 - 2\big].
\]
As ra\'izes s\~ao:
\[
\lambda_1 = 2, \quad \lambda_2 = 2+\sqrt{2}, \quad \lambda_3 = 2-\sqrt{2}.
\]
Como todos $\lambda_i > 0$, conclu\'imos que $A$ \'e \textbf{positiva definida}.

\subsection*{3) Interpreta\c c\~ao e Rela\c c\~ao com a Forma Quadr\'atica}

A forma quadr\'atica associada a uma matriz sim\'etrica $A$ \'e
\[
q_A(x) = x^\top A x.
\]
Ela mede como a transforma\c c\~ao linear $A$ atua sobre o vetor $x$, combinando sua dire\c c\~ao e magnitude.  
A an\'alise de $q_A(x)$ \'e central na classifica\c c\~ao de pontos cr\'iticos em otimiza\c c\~ao e na determina\c c\~ao da curvatura local de superf\'icies.

\paragraph{Decomposi\c c\~ao espectral.}
Se $A$ \'e sim\'etrica, existe uma matriz ortogonal $Q$ e uma matriz diagonal $\Lambda = \mathrm{diag}(\lambda_1, \dots, \lambda_n)$ tais que:
\[
A = Q \Lambda Q^\top, \quad \text{com } Q^\top Q = I.
\]
Substituindo em $q_A(x)$:
\[
x^\top A x = x^\top Q \Lambda Q^\top x.
\]
Definindo $y = Q^\top x$ (ou seja, uma mudan\c ca de base para o sistema de autovetores de $A$), temos:
\[
q_A(x) = y^\top \Lambda y = \sum_{i=1}^n \lambda_i y_i^2.
\]

\paragraph{Significado.}
Essa express\~ao mostra que, no sistema de coordenadas formado pelos autovetores de $A$, a forma quadr\'atica se decomp\~oe em uma soma de termos independentes, cada um ponderado por um autovalor $\lambda_i$:
\[
q_A(x) = \lambda_1 y_1^2 + \lambda_2 y_2^2 + \cdots + \lambda_n y_n^2.
\]
Assim, o sinal de $q_A(x)$ depende exclusivamente dos sinais dos $\lambda_i$, pois $y_i^2 \ge 0$.

\paragraph{Consequ\^encia imediata.}
\begin{itemize}
  \item Se todos $\lambda_i > 0$, ent\~ao $q_A(x) > 0$ para todo $x \ne 0$ $\Rightarrow$ $A$ \'e \textbf{positiva definida}.
  \item Se todos $\lambda_i < 0$, ent\~ao $q_A(x) < 0$ para todo $x \ne 0$ $\Rightarrow$ $A$ \'e \textbf{negativa definida}.
  \item Se alguns $\lambda_i > 0$ e outros $< 0$, o sinal de $q_A(x)$ varia conforme a dire\c c\~ao $\Rightarrow$ $A$ \'e \textbf{indefinida}.
  \item Se todos $\lambda_i \ge 0$ mas algum \'e zero, ent\~ao $q_A(x)\ge0$ $\Rightarrow$ $A$ \'e \textbf{positiva semidefinida}.
\end{itemize}

\paragraph{Intui\c c\~ao geom\'etrica.}
Cada autovetor $v_i$ define um \textbf{eixo principal} da superf\'icie descrita por $z = q_A(x)$.  
O autovalor correspondente $\lambda_i$ controla a \textbf{curvatura} (ou “abertura”) da par\'abola nessa dire\c c\~ao:
\begin{itemize}
  \item $\lambda_i > 0$: a superf\'icie curva para cima — dire\c c\~ao de m\'inimo (vale de paraboloide).
  \item $\lambda_i < 0$: a superf\'icie curva para baixo — dire\c c\~ao de m\'aximo (c\'upula invertida).
  \item mistura de sinais: a superf\'icie tem \textbf{curvatura de sela}, positiva em algumas dire\c c\~oes e negativa em outras.
\end{itemize}

\paragraph{Conex\~ao com otimiza\c c\~ao.}
No contexto da an\'alise de uma fun\c c\~ao $f(x)$ em um ponto cr\'itico $x^\ast$, a Hessiana $H = \nabla^2 f(x^\ast)$ desempenha o papel de $A$.  
\begin{itemize}
  \item Se $H$ \'e positiva definida, $f$ tem um m\'inimo local estrito em $x^\ast$.
  \item Se $H$ \'e negativa definida, $f$ tem um m\'aximo local estrito.
  \item Se $H$ \'e indefinida, $x^\ast$ \'e um ponto de sela.
\end{itemize}
Assim, a classifica\c c\~ao de matrizes via autovalores fornece uma ponte direta entre \textbf{\'algebra linear} e \textbf{otimiza\c c\~ao n\~ao linear}.


\subsection{Crit\'erio de Sylvester (Menores Principais)}

\subsection*{1) Motiva\c c\~ao}
Embora o crit\'erio dos autovalores seja o mais direto numericamente, ele nem sempre \'e conveniente em demonstra\c c\~oes anal\'iticas.  
O \textbf{Crit\'erio de Sylvester} fornece uma maneira puramente alg\'ebrica de verificar a definitude de uma matriz sim\'etrica $A$, analisando os \textbf{determinantes dos menores principais}.

Esse m\'etodo evita o c\'alculo expl\'icito dos autovalores e \'e amplamente usado em \textbf{an\'alise te\'orica}, \textbf{provas formais} e \textbf{verifica\c c\~ao simb\'olica} de positividade.

\subsection*{2) Defini\c c\~ao Formal}
Seja $A \in \mathbb{R}^{n\times n}$ uma matriz sim\'etrica.  
Definimos seus \textbf{menores principais} $A_k$ como as submatrizes $k\times k$ obtidas pelos primeiros $k$ elementos das linhas e colunas:
\[
A_k = 
\begin{bmatrix}
a_{11} & \dots & a_{1k} \\
\vdots & \ddots & \vdots \\
a_{k1} & \dots & a_{kk}
\end{bmatrix}.
\]
O determinante de cada $A_k$, denotado $\Delta_k = \det(A_k)$, \'e chamado de \textbf{menor principal de ordem $k$}.

O Crit\'erio de Sylvester afirma:

\begin{center}
\begin{tabular}{ll}
$A$ positiva definida (PD) & $\iff$ $\Delta_k > 0$ para todo $k=1,\dots,n$.\\[3pt]
$A$ negativa definida (ND) & $\iff$ os sinais de $\Delta_k$ alternam: $\Delta_1 < 0$, $\Delta_2 > 0$, $\Delta_3 < 0$, etc.
\end{tabular}
\end{center}

\paragraph{Casos semidefinidos:}
\begin{itemize}
  \item Se todos $\Delta_k \ge 0$, $A$ pode ser \textbf{positiva semidefinida}, mas n\~ao necessariamente (o crit\'erio n\~ao \'e suficiente nesse caso).  
  \item Analogamente, $\Delta_k$ alternando em $\le 0, \ge 0, \le 0, \dots$ pode indicar \textbf{negativa semidefinida}.
\end{itemize}

\subsection*{3) Exemplo Num\'erico}
Considere:
\[
A =
\begin{bmatrix}
2 & -1 & 0 \\
-1 & 2 & -1 \\
0 & -1 & 2
\end{bmatrix}.
\]
Calculamos os menores principais:
\[
\begin{aligned}
A_1 &= [2], & \quad \Delta_1 &= 2,\\
A_2 &= 
\begin{bmatrix}
2 & -1\\
-1 & 2
\end{bmatrix}, & \quad \Delta_2 &= (2)(2) - (-1)^2 = 3,\\
A_3 &= A, & \quad \Delta_3 &= \det(A) = 4.
\end{aligned}
\]
Todos os determinantes s\~ao positivos:
\[
\Delta_1 = 2 > 0, \quad \Delta_2 = 3 > 0, \quad \Delta_3 = 4 > 0.
\]
Portanto, pelo Crit\'erio de Sylvester, $A$ \'e \textbf{positiva definida}.  
Esse resultado coincide com a classifica\c c\~ao obtida via autovalores.

\subsection*{4) Por que o Crit\'erio de Sylvester funciona}

A ideia central vem da decomposi\c c\~ao de Cholesky:  
uma matriz sim\'etrica $A$ \'e positiva definida se, e somente se, ela pode ser escrita como
\[
A = L L^\top,
\]
onde $L$ \'e triangular inferior com diagonais positivas.

Durante o processo de fatora\c c\~ao, cada $\Delta_k = \det(A_k)$ representa o produto dos quadrados dos primeiros $k$ elementos diagonais de $L$:
\[
\Delta_k = (l_{11} l_{22} \cdots l_{kk})^2.
\]
Portanto:
\[
\Delta_k > 0 \ \forall k \quad \Longleftrightarrow \quad l_{ii} > 0 \ \forall i \quad \Longleftrightarrow \quad A = L L^\top \text{ com } L \text{ n\~ao singular.}
\]
Assim, o sinal dos menores principais reflete diretamente a \textbf{positividade das curvaturas parciais} da forma quadr\'atica $x^\top A x$, medida ao longo das primeiras $k$ dire\c c\~oes coordenadas.

\paragraph{Interpreta\c c\~ao geom\'etrica:}
Cada $\Delta_k$ corresponde ao “volume” (ou \textit{determinante}) da elips\'oide definida pela restri\c c\~ao de $A$ ao subespa\c co gerado pelas primeiras $k$ vari\'aveis.  
Se algum $\Delta_k \le 0$, significa que a forma quadr\'atica n\~ao \'e positiva em todas as dire\c c\~oes desse subespa\c co — indicando perda de convexidade parcial.

\subsection*{5) Comparativo com o crit\'erio dos autovalores}
\begin{itemize}
  \item Ambos caracterizam positividade, mas o de Sylvester evita c\'alculo espectral.
  \item O de autovalores \'e mais direto computacionalmente (usado em pr\'atica).
  \item O de Sylvester \'e mais pr\'atico em demonstra\c c\~oes anal\'iticas e provas simb\'olicas.
\end{itemize}

\subsection*{S\'intese dos Crit\'erios de Classifica\c c\~ao}

A Tabela~\ref{tab:comparacao-criterios} resume as principais caracter\'isticas dos tr\^es crit\'erios mais utilizados para determinar a definitude de uma matriz sim\'etrica: o dos autovalores, o de Sylvester e o baseado na fatora\c c\~ao de Cholesky.

\begin{table}[h]
\centering
\small
\setlength{\tabcolsep}{5pt}
\renewcommand{\arraystretch}{1.2}
\begin{tabularx}{\textwidth}{@{}lXXXX@{}}
\toprule
\textbf{Crit\'erio} &
\textbf{Base te\'orica} &
\textbf{Informa\c c\~ao usada} &
\textbf{Vantagens} &
\textbf{Limita\c c\~oes} \\
\midrule
\textbf{Autovalores} &
Decomposi\c c\~ao espectral $A = Q\Lambda Q^\top$ &
Sinais dos autovalores $\lambda_i$ &
Interpreta\c c\~ao geom\'etrica direta; base para classifica\c c\~ao via Hessiana &
Requer c\'alculo espectral; mais pesado simbolicamente \\

\textbf{Sylvester} &
Menores principais $\det(A_k)$ &
Determinantes dos subblocos superiores &
Alg\'ebrico e simb\'olico; ideal em provas anal\'iticas &
N\~ao conclusivo em casos semidefinidos \\

\textbf{Cholesky} &
Fatora\c c\~ao $A = L L^\top$ &
Elementos diagonais de $L$ &
Computacionalmente eficiente; confirma positividade direta &
Apenas aplic\'avel se $A$ for positiva definida \\
\bottomrule
\end{tabularx}
\caption{Comparativo entre os crit\'erios de classifica\c c\~ao de matrizes sim\'etricas.}
\label{tab:comparacao-criterios}
\end{table}

\paragraph{Resumo.}
Os tr\^es crit\'erios s\~ao equivalentes para matrizes positivas definidas, mas diferem em prop\'osito:
\begin{itemize}
  \item \textbf{Autovalores:} fornece interpreta\c c\~ao geom\'etrica e liga\c c\~ao direta com curvatura e Hessianas;
  \item \textbf{Sylvester:} \'util para demonstra\c c\~oes formais e verifica\c c\~ao simb\'olica;
  \item \textbf{Cholesky:} preferido em implementa\c c\~oes num\'ericas e verifica\c c\~ao computacional.
\end{itemize}


\section{Exemplo Num\'erico (3$\times$3)}
Considere
\[
A=\begin{bmatrix}
4 & 1 & -1\\
1 & 3 & 0\\
-1& 0 & 2
\end{bmatrix}.
\]
Podemos (i) computar autovalores para classificar, (ii) visualizar $q_A(x)=x^\top A x$ em $\mathbb{R}^2$ fixando $x=(u,v,0)$, e (iii) aplicar Sylvester.

\section{Exemplo pr\'atico em Python (autovalores, Sylvester e gr\'afico 3D)}
\noindent
O c\'odigo abaixo:
\begin{enumerate}
  \item classifica $A$ via autovalores com toler\^ancia num\'erica;
  \item reporta menores principais (Sylvester);
  \item plota $z = [u,v,0]^\top A [u,v,0]$ em uma malha 2D (gr\'afico 3D) para intui\c c\~ao geom\'etrica.
\end{enumerate}

\begin{lstlisting}
import numpy as np
import matplotlib.pyplot as plt
from mpl_toolkits.mplot3d import Axes3D  # noqa: F401 (ensures 3D projection is registered)

# ---- Matrix (3x3) ----
A = np.array([[4.0, 1.0, -1.0],
              [1.0, 3.0,  0.0],
              [-1.0, 0.0, 2.0]])

# ---- Classification via eigenvalues ----
def classify_definiteness(A, tol=1e-10):
    lam = np.linalg.eigvalsh(A)  # symmetric eigenvalues
    pos = np.all(lam >  tol)
    nonneg = np.all(lam >= -tol)
    neg = np.all(lam < -tol)
    nonpos = np.all(lam <=  tol)
    if pos:
        return "Positiva definida (PD)", lam
    if nonneg and not pos:
        return "Positiva semidefinida (PSD)", lam
    if neg:
        return "Negativa definida (ND)", lam
    if nonpos and not neg:
        return "Negativa semidefinida (NSD)", lam
    return "Indefinida", lam

label, eigvals = classify_definiteness(A)
print("Classificacao:", label)
print("Autovalores  :", eigvals)

# ---- Sylvester: leading principal minors ----
def leading_principal_minors(A):
    dets = []
    for k in range(1, A.shape[0]+1):
        dets.append(np.linalg.det(A[:k, :k]))
    return dets

dets = leading_principal_minors(A)
print("Menores principais (det A_k):", dets)

# ---- 3D plot of z = x^T A x with x=(u,v,0) ----
u = np.linspace(-2.5, 2.5, 200)
v = np.linspace(-2.5, 2.5, 200)
U, V = np.meshgrid(u, v)

# x = [u, v, 0]; z = [u,v,0]^T A [u,v,0]
Z = (A[0,0]*U**2 + 2*A[0,1]*U*V + A[1,1]*V**2)  # terms with the third comp zero
# Cross terms with third coord vanish because x3=0

fig = plt.figure(figsize=(7,5))
ax = fig.add_subplot(111, projection='3d')
ax.plot_surface(U, V, Z, alpha=0.7, linewidth=0, antialiased=True)
ax.set_xlabel('u')
ax.set_ylabel('v')
ax.set_zlabel('z = x^T A x (x3=0)')
ax.set_title(f"Forma quadratica (x3=0) --- {label}")
plt.tight_layout()
plt.show()
\end{lstlisting}

\paragraph{Observa\c c\~ao.}
Para visualizar cortes diferentes, altere a proje\c c\~ao (ex.: $x=(u,0,w)$ ou $x=(0,v,w)$) trocando os \'indices correspondentes na express\~ao de $Z$.


\section{Resumo (tabela de bolso)}

\begin{table}[!ht]
\centering
\small
\setlength{\tabcolsep}{5pt}
\renewcommand{\arraystretch}{1.15}
\begin{tabularx}{\textwidth}{@{}lccX@{}}
\toprule
\textbf{Classe} & \textbf{Sinal de $x^\top A x$} & \textbf{Autovalores (sim\'etrica)} & \textbf{Crit\'erio de Sylvester (ment. principais)} \\
\midrule
PD & $>0$ para $x\ne 0$ & todos $>0$ & $\det(A_k)>0,\ \forall k$ \\
PSD & $\ge 0$ & todos $\ge 0$ & n\~ao conclusivo isoladamente (ver autovalores) \\
ND & $<0$ para $x\ne 0$ & todos $<0$ & sinais alternados: $-,+,-,\dots$ \\
NSD & $\le 0$ & todos $\le 0$ & n\~ao conclusivo isoladamente (ver autovalores) \\
Indefinida & sinais opostos & autovalores de sinais mistos & n/a \\
\bottomrule
\end{tabularx}
\caption{Classifica\c c\~ao de matrizes sim\'etricas via forma quadr\'atica.}
\end{table}

\section{Dicas Pr\'aticas}
\begin{itemize}
  \item Em c\'alculo computacional, use toler\^ancia ($\sim 10^{-10}$) ao decidir sinais de autovalores.
  \item Para Hessianas, prefira autovalores (via \texttt{eigvalsh}); Sylvester \'e \'util para prova anal\'itica.
  \item Em otimiza\c c\~ao, a definitude da Hessiana no \'otimo influencia a taxa de converg\^encia (Newton/quasi-Newton).
\end{itemize}

\end{document}
